%-------------------------------------------------------------------------------
% Resume
%-------------------------------------------------------------------------------
\documentclass[a4paper,10pt]{article}

\usepackage{url}
\usepackage{parskip}
\usepackage[dvipsnames]{xcolor}
\usepackage[scale=0.91]{geometry}
\usepackage{tabularx}
\usepackage{titlesec}
\usepackage{fontawesome5}
\usepackage[unicode,draft=false]{hyperref}
\definecolor{linkcolour}{rgb}{0,0.2,0.6}
\definecolor{experienceblue}{HTML}{2F5D8A}
\hypersetup{colorlinks,breaklinks,urlcolor=linkcolour,linkcolor=linkcolour}

\titleformat{\section}{\large\scshape\raggedright}{}{0em}{}[\titlerule]
\titlespacing{\section}{0pt}{7pt}{5pt}
\pagestyle{empty}

\newcommand{\datedentry}[2]{%
\begin{tabularx}{\linewidth}{@{}X r@{}}
#1 & #2
\end{tabularx}}

\newenvironment{researchentry}[2]{%
\datedentry{\textbf{#1}}{#2}
\vspace{-2pt}
\par\noindent\ignorespaces
}{%
\par
}

\begin{document}

%-------------------------------------------------------------------------------
% Header
%-------------------------------------------------------------------------------
\begin{center}
    {\LARGE \textbf{Zixun Huang}}\\[3pt]
    {\small
    \href{mailto:zixunh@wharton.upenn.edu}{\faEnvelope\ zixunh@wharton.upenn.edu} \; | \;
    \href{https://alexhuang13.github.io/}{\faGlobe\ Homepage} \; | \;
    \href{https://github.com/alexhuang13}{\faGithub\ GitHub} \; | \;
    \href{https://scholar.google.com/citations?user=n7NzBzAAAAAJ}{\faGraduationCap\ Google Scholar} \; | \;
    \href{https://www.linkedin.com/in/zixun-huang-b80ab4351/}{\faLinkedin\ LinkedIn}}
\end{center}
\vspace{-9pt}

%-------------------------------------------------------------------------------
\section{Research Interests}
Theoretical foundations of machine learning, with an emphasis on \textbf{Optimization Theory \& Scaling Laws}, \textbf{LLM Reasoning \& Post-Training \& Agentic RL}, and \textbf{applications of AI}.

%-------------------------------------------------------------------------------
\section{Education}
\datedentry{\textbf{University of Pennsylvania, Wharton School}}{Aug. 2026 -- Present}
\vspace{-2pt}
\hspace{1em}\textit{Ph.D. in Statistics and Data Science}

\datedentry{\textbf{Peking University, School of Mathematical Sciences}}{Sept. 2022 -- Jun. 2026}
\vspace{-2pt}
\hspace{1em}\textit{B.S. in Statistics}; Elite Undergraduate Program for Applied Mathematics; GPA: 3.7/4.0

\datedentry{\textbf{University of California, Berkeley}}{Jan. 2025 -- Aug. 2025}
\vspace{-2pt}
\hspace{1em}\textit{Visiting Student}; GPA: 4.0/4.0

%-------------------------------------------------------------------------------
\section{Publications}
[1] Zhongzhi Li, Yucheng Shi, Zongxia Li, Ruhan Wang, Anhao Li, \textbf{Zixun Huang}, Junyao Yang, Lei Ke, Ninghao Liu, Haitao Mi, Leowei Liang. \textit{Recursive Synthesis for Long-Horizon Terminal Tasks}. \textbf{Preprint}. \href{https://arxiv.org/abs/2608.05466}{[Paper]} \href{https://huggingface.co/collections/Zhongzhi1228/recursive-task-synthesis}{[Hugging Face]} \href{https://zhongzhi660.github.io/recursive-verified-synthesis-site/?case=jobs-diff-01-3341b098}{[Website]}

[2] \textbf{Zixun Huang*}, Jiayi Sheng*, Zeyu Zheng. \textit{Variance-Aware Baselines and Adaptive Learning Rates for Reinforcement Learning with Verifiable Rewards}. \textbf{Under review}. \href{https://arxiv.org/abs/2511.23310}{[Paper]}

[3] Yuhao Sun, Zekun Wu, \textbf{Zixun Huang}, Peijie Zhou. \textit{TracingFlow: A Simulation-Free Trajectory Inference Framework Based on Second-Order Dynamics}. \textbf{Under review}.

[4] Binghui Li*, Fengling Chen*, \textbf{Zixun Huang*}, Lean Wang*, Lei Wu. \textit{Functional Scaling Laws in Kernel Regression: Loss Dynamics and Learning Rate Schedules}. \textcolor[HTML]{8B0000}{\textbf{NeurIPS 2025 Spotlight}}. \href{https://arxiv.org/abs/2509.19189}{[Paper]} {\footnotesize * Equal contribution.}

%-------------------------------------------------------------------------------
\section{Experience}
\noindent\textsf{Optimization Theory \& Scaling Laws}\par
\vspace{-2pt}
\datedentry{\faCaretRight\ \textcolor{experienceblue}{\textbf{Peking University}}, \textit{Research Intern}, Beijing}{May 2023 -- Aug. 2026}
\vspace{-3pt}
\hspace{1em}\textit{Supervisor: Assistant Professor Lei Wu}

\noindent\textsf{LLM Reasoning \& Post-Training \& Agentic RL}\par
\vspace{-2pt}
\datedentry{\faCaretRight\ \textcolor{experienceblue}{\textbf{Tencent Hy LLM -- Qingyun Talent Program}}, \textit{Research Intern}, Beijing}{Apr. 2026 -- Present}
\vspace{-3pt}
\hspace{1em}\textit{Supervisor: Senior Researcher Kishan Panaganti}

\datedentry{\faCaretRight\ \textcolor{experienceblue}{\textbf{University of California, Berkeley}}, \textit{Research Intern}, Berkeley, CA}{Jan. 2025 -- Aug. 2025}
\vspace{-3pt}
\hspace{1em}\textit{Supervisor: Associate Professor Zeyu Zheng}

%-------------------------------------------------------------------------------
\section{Selected Honors and Awards}
\begin{tabularx}{\linewidth}{@{}X r@{}}
Grand Prize, ``Challenge Cup'' Competition & 2026 \\
Weiming Scholar, Peking University & 2026 \\
Hong Sheng Scholarship & 2024 \\
First Prize, Chinese Mathematics Contest & 2024 \\
Yau Contest Groups Prize & 2024 \\
Xiaomi Scholarship & 2023 \\
Gold Prize, Chinese Mathematical Olympiad & 2021 \\
Silver Prize, Chinese Mathematical Olympiad & 2020
\end{tabularx}

%-------------------------------------------------------------------------------
\section{Skills}
\begin{tabularx}{\linewidth}{@{}l X@{}}
\textbf{Programming:} & Python (NumPy, pandas, scikit-learn, PyTorch), C++, MATLAB, \LaTeX \\
\textbf{Language:} & GRE V160/Q170; GRE Math 970 (95th percentile); TOEFL 101
\end{tabularx}

\end{document}
